\scanpage{206}

Soit $f:X\longrightarrow S$ un morphisme de schémas, $f^{\mathrm{pl}}$ le morphisme correspondant des sites fppf.

\textcircled{1} $P^c=P^c_{X/S}:\ (\mathrm{Sch})_{/S}^{\unsure{o|\circ}}\longrightarrow(\mathrm{Ab})$, \deleted{défini par}\inserted{faisceau associé au préfaisceau} \deleted{$P^c(S')=$} $S'\longmapsto$ classes de modules inversibles sur $X_{S'}$, munis d'une connexion rel/$S'$

\[
0\longrightarrow\frac{f_*^{\mathrm{pl}}(\underline\Omega^1_{X/S})}{\mathrm{Im}\, f_*^{\mathrm{pl}}(\underline O^*_{X/S})}
\longrightarrow P^c_{X/S}\longrightarrow\underline{\mathrm{Pic}}_{X/S}
\longrightarrow\mathrm R^1 f_*^{\mathrm{pl}}(\underline\Omega^1_{X/S})\tag*{(1.1)}
\]

\margin{$f_*^{\mathrm{pl}}\bigl(\underline O^*_{X/S}\bigr)\longrightarrow f_*^{\mathrm{pl}}(\Omega^1_{X/S})$ donné par $\varphi\longmapsto d\varphi/\varphi$.}

% slanted

\margin{Cet hom est nul si $f$ est propre et séparable, \inserted{de prés. finie} p.ex. propre et lisse}

\uline{NB} Si $f$ \deleted{propre et \unsure{gén. sép.|plat}}\inserted{cohérent}, et $\underline\Omega^1_{X/S}$ cohom. plat/$S$ en dim $0$, ou $\mathrm R^1 f_*(\underline\Omega^1_{X/S})$ commute au chgt de base, alors pour tt schéma abélien $A$ sur $S$, tout hom $B\mathbin{\unsure{-|\longrightarrow}}\mathrm R^1 f_*^{\mathrm{pl}}(\underline\Omega^{1}_{X/S})$ est nul. donc si on a un hom $B\overset{u}{\longrightarrow}\underline{\mathrm{Pic}}_{X/S}$, on a $u(B)\subset\mathrm{Im}(P^c\longrightarrow\underline{\mathrm{Pic}}_{X/S})$, \unsure{d'où|donc} suite exacte

\[
\cdot\mathbin{-}f_*^{\mathrm{pl}}(\Omega^1_{X/S})\longrightarrow u^*(P^c)\longrightarrow B\longrightarrow1
\tag*{\unsure{(1.2')|(1.2)}}
\]

Ceci s'applique en particulier si \uline{$f$} propre sur $S$, \deleted{\unsure{do}} donc $\underline{\mathrm{Pic}}^0_{X/S}$ est défini, et si $\boxed{B=\underline{\mathrm{Pic}}^0_{X/S}}$ est un schéma abélien. Alors, $u$ étant l'inclusion, on aura $u^*(P^c)=(P^c)^0$, et on trouve suite exacte

\[
\cdot\longrightarrow f_*^{\mathrm{pl}}(\underline\Omega^1_{X/S})\longrightarrow(P^c_{X/S})^0
\longrightarrow\underset{\substack{\Vert\\\underline{\mathrm{Pic}}^0_{X/S}}}{B}\longrightarrow\cdot
\tag*{\unsure{(1.2)|(1.2')|(1.3)}}
\]

\uline{Re NB} \unsure{$\alpha_1$|ceci|ce}, et s'il est défini, \unsure{$\alpha_2$|ceci|ce|cas}, est fonctoriel en $X/S$, et commute au chgt de base.

\textcircled{2} $P^{ci}=P^{ci}_{X/S}$, sous-Groupe de $P^c_{X/S}$, faisceau associé au préfaisceau $S'\longmapsto$ classes de modules inversibles sur $X_{S'}$, munis d'une connexion rel. à $S'$ qui est \uline{à courbure nulle}

% archive sheet number: 205
